% ============================================================
% maw_comun.tex --- preambulo compartido MAW Soluciones
% Lo usan PROPUESTA, CONTRATO y RESUMEN_EJECUTIVO.
% El documento debe definir ANTES del \input:
%   \docTipo, \docPrefijo, \fechaProp, \tarifa, \demoURL,
%   \clienteCorto, \clienteLargo
% Compilar con xelatex (Montserrat via fontspec), dos pasadas.
% ============================================================

% ============================================================
% PAQUETES
% ============================================================
\usepackage[spanish,es-tabla,es-noindentfirst]{babel}
\usepackage{fontspec}
\usepackage{montserrat}
\renewcommand*\familydefault{\sfdefault}
\usepackage[T1]{fontenc}
\usepackage{geometry}
\usepackage{graphicx}
\usepackage[table]{xcolor}
\usepackage{tabularx}
\usepackage{array}
\usepackage{booktabs}
\usepackage{ragged2e}
\usepackage{enumitem}
\usepackage{tcolorbox}
\usepackage{needspace}
\usepackage{fancyhdr}
\usepackage{siunitx}
\usepackage{hyperref}
\usepackage{tikz}
\usetikzlibrary{positioning, arrows.meta, shapes.geometric, calc}

\geometry{top=2cm, bottom=2cm, left=1.8cm, right=1.8cm}
\setlength{\headheight}{26pt}
\addtolength{\topmargin}{-14pt}

\sisetup{group-separator={,}, group-minimum-digits=4, output-decimal-marker={.}}

% ============================================================
% COLORES --- MAW Soluciones / PharmaWeigh
% (los mismos tokens que usa la app: src/app/globals.css)
% ============================================================
\definecolor{azulcom}{HTML}{1E40AF}    % acento principal (primary)
\definecolor{doradocom}{HTML}{2563EB}   % filete y acento secundario (blue-600)
\definecolor{grisbar}{HTML}{1E293B}     % slate-800: barra lateral de la app
\definecolor{grisfondo}{HTML}{F1F5F9}   % slate-100: fondo de tablas
\definecolor{gristexto}{HTML}{64748B}   % slate-500: texto secundario
\definecolor{verdeok}{HTML}{16A34A}     % éxito / dentro de tolerancia
\definecolor{ambarpend}{HTML}{D97706}   % advertencia / pendiente
\definecolor{rojonota}{HTML}{DC2626}    % error / fuera de tolerancia
\definecolor{tinta}{HTML}{0F172A}       % texto principal

\hypersetup{colorlinks=true, linkcolor=azulcom, urlcolor=doradocom}

% ============================================================
% ENCABEZADO / PIE
% ============================================================
\pagestyle{fancy}
\fancyhf{}
\fancyhead[L]{\footnotesize\color{gristexto}\textbf{MAW Soluciones} \textbar{} \docTipo}
\fancyhead[R]{\footnotesize\color{gristexto}\fechaProp}
\fancyfoot[C]{\footnotesize\color{gristexto}\thepage}
\renewcommand{\headrulewidth}{0.6pt}
\renewcommand{\headrule}{\hbox to\headwidth{\color{azulcom}\leaders\hrule height \headrulewidth\hfill}}

% ============================================================
% TITULOS Y CAJAS
% ============================================================
\newcommand{\tituloCot}[1]{%
    \noindent{\LARGE\bfseries\color{azulcom}\docPrefijo{} \textbar{} #1}\par
    \vspace{2pt}
    {\color{doradocom}\rule{\linewidth}{1.8pt}}\par
    \vspace{14pt}
}

\newcommand{\seccion}[1]{%
    {\large\bfseries\color{azulcom}#1}\\[4pt]
    {\color{doradocom}\rule{\linewidth}{1pt}}\par
    \vspace{8pt}
}

\newcommand{\subseccion}[1]{%
    \vspace{4pt}
    \noindent{\bfseries\color{azulcom}#1}\par
    \vspace{4pt}
}

\newtcolorbox{cajaNota}{
    colback=grisfondo, colframe=gristexto!60, rounded corners,
    boxrule=1pt, left=10pt, right=10pt, top=6pt, bottom=6pt
}

\newtcolorbox{cajaInfo}[1][]{
    colback=azulcom!5, colframe=azulcom,
    fonttitle=\bfseries\color{white}, colbacktitle=azulcom, title=#1,
    rounded corners, boxrule=1pt, left=8pt, right=8pt, top=6pt, bottom=6pt
}

\newtcolorbox{cajaTotal}{
    colback=verdeok!8, colframe=verdeok, rounded corners,
    boxrule=1.5pt, left=10pt, right=10pt, top=10pt, bottom=10pt
}

\newtcolorbox{cajaAviso}[1][]{
    colback=rojonota!6, colframe=rojonota,
    fonttitle=\bfseries\color{white}, colbacktitle=rojonota, title=#1,
    rounded corners, boxrule=1.2pt, left=8pt, right=8pt, top=6pt, bottom=6pt
}

% ============================================================
% PENDIENTES --- nada que el cliente o Fernando no hayan
% confirmado se escribe como cifra. Se marca y se ve.
% ============================================================
\newtcolorbox{cajaPendiente}[1][Pendiente de autorización]{
    colback=ambarpend!8, colframe=ambarpend,
    fonttitle=\bfseries\color{white}, colbacktitle=ambarpend, title=#1,
    rounded corners, boxrule=1.2pt, left=8pt, right=8pt, top=6pt, bottom=6pt
}

% \PENDIENTE{quien}{que} --- marca visible en linea
\newcommand{\PENDIENTE}[2]{%
    {\setlength{\fboxsep}{2.5pt}%
     \colorbox{ambarpend!20}{\color{ambarpend!80!black}\footnotesize\bfseries PENDIENTE (#1): #2}}%
}

% Version compacta para celdas de tabla
\newcommand{\PENDmini}{%
    {\setlength{\fboxsep}{2pt}%
     \colorbox{ambarpend!20}{\color{ambarpend!80!black}\scriptsize\bfseries pend.}}%
}

% ============================================================
% MOTOR DE CIFRAS
% ------------------------------------------------------------
% Mientras \cifrasDefinidas sea falso, toda hora e importe se
% imprime como PENDIENTE. Para activar el documento:
%   1. Fernando autoriza tarifa y horas.
%   2. Se pone \cifrasDefinidastrue y se llenan \tarifaNum y
%      los \def de horas de cada renglon.
%   3. Los subtotales, el total, la mensualidad y el punto de
%      equilibrio se calculan solos con estas macros.
% ============================================================
\newif\ifcifrasDefinidas
\cifrasDefinidasfalse

\providecommand{\tarifaNum}{0}     % MXN por hora de ingenieria (sin separadores)

% \HR{expresion} --- horas (acepta sumas: \HR{\hAa+\hAb})
\newcommand{\HR}[1]{\ifcifrasDefinidas\fpeval{#1}\else\PENDmini\fi}

% \IM{expresion} --- importe = horas x tarifa
\newcommand{\IM}[1]{\ifcifrasDefinidas\$\num{\fpeval{round((#1)*\tarifaNum,0)}}\else\PENDmini\fi}

% \MXN{numero} --- importe literal ya autorizado (mensualidad, infra)
\newcommand{\MXN}[1]{\ifcifrasDefinidas\$\num{#1}\else\PENDmini\fi}

% \RAZON{a}{b} --- cociente redondeado hacia arriba (punto de equilibrio)
\newcommand{\RAZON}[2]{\ifcifrasDefinidas\num{\fpeval{ceil((#1)/(#2))}}\else\PENDmini\fi}

% ============================================================
% TABLAS
% ============================================================
\newcolumntype{R}[1]{>{\RaggedLeft\arraybackslash}p{#1}}
\newcolumntype{C}[1]{>{\centering\arraybackslash}p{#1}}
\newcolumntype{L}[1]{>{\RaggedRight\arraybackslash}p{#1}}

\renewcommand{\arraystretch}{1.3}
